\documentclass[11pt,letterpaper]{article}
\usepackage[T1]{fontenc}
\usepackage[utf8]{inputenc}
\usepackage{lmodern}
\usepackage[margin=0.95in]{geometry}
\usepackage{amsmath,amssymb,amsthm}
\usepackage{enumitem}
\usepackage{microtype}
\usepackage{xcolor}
\usepackage{booktabs,array}
\usepackage{pict2e}
\usepackage{fancyhdr}
\usepackage{hyperref}
\definecolor{mathjockgreen}{HTML}{355439}
\hypersetup{colorlinks=true,linkcolor=mathjockgreen,urlcolor=mathjockgreen,
  pdftitle={Fraleigh A First Course in Abstract Algebra: worked solutions},
  pdfauthor={Mathjock},pdfsubject={Detailed worked solutions, seventh edition}}
\setlength{\parindent}{0pt}
\setlength{\parskip}{0.55em}
\setlength{\headheight}{14pt}
\setlist[enumerate]{leftmargin=*,itemsep=0.4em,topsep=0.4em}
\setlist[itemize]{leftmargin=*,itemsep=0.25em}
\allowdisplaybreaks
\emergencystretch=2em
% Leave a visible gap after two-digit exercise numbers in the contents.
\makeatletter
% Chapter.section and supplementary labels need more room than article's
% default 1.5em section-number box, especially for labels such as 3.S.
\renewcommand*{\l@section}[2]{\begingroup\bfseries
  \@dottedtocline{1}{0em}{3em}{#1}{#2}\endgroup}
\renewcommand*{\l@subsection}{\@dottedtocline{2}{1.5em}{3.5em}}
\makeatother
\pagestyle{fancy}
\fancyhf{}
\fancyhead[L]{\small Mathjock \textbar\ Fraleigh, seventh edition}
\fancyhead[R]{\small Step-by-step solutions}
\fancyfoot[C]{\thepage}
\newcommand{\N}{\mathbb{N}}
\newcommand{\Z}{\mathbb{Z}}
\newcommand{\Q}{\mathbb{Q}}
\newcommand{\R}{\mathbb{R}}
\newcommand{\abs}[1]{\left|#1\right|}
\makeatletter
\newcommand{\exercise}[2]{\subsection*{Exercise #1: #2}\addcontentsline{toc}{subsection}{Exercise #1: #2}\def\@currentlabel{#1}\label{ex:#1}}
\makeatother
\newcommand{\goal}{\par\textbf{What we must show.} }
\newcommand{\idea}{\par\textbf{Plan.} }
\newcommand{\step}[1]{\par\medskip\textbf{#1}\par}
\newcommand{\answer}[1]{\par\textbf{Conclusion.} #1\par}
\newcommand{\C}{\mathbb{C}}
\DeclareMathOperator{\Span}{Span}
\DeclareMathOperator{\Col}{Col}
\DeclareMathOperator{\Nul}{Nul}
\DeclareMathOperator{\rank}{rank}
\DeclareMathOperator{\dist}{dist}
\newcommand{\practice}[2]{\subsection*{Practice problem #1: #2}\addcontentsline{toc}{subsection}{Practice problem #1: #2}}

\DeclareMathOperator{\Aut}{Aut}
\DeclareMathOperator{\Inn}{Inn}
\DeclareMathOperator{\im}{im}
\DeclareMathOperator{\Ker}{Ker}
\DeclareMathOperator{\ord}{ord}
\DeclareMathOperator{\Gal}{Gal}
\DeclareMathOperator{\lcm}{lcm}
\DeclareMathOperator{\Hom}{Hom}
\newcommand{\gen}[1]{\langle #1\rangle}
\newcommand{\bookinfo}{%
\textbf{Text and numbering.} John B. Fraleigh, \emph{A First Course in Abstract
Algebra}, seventh edition (2003), following the scanned textbook linked on
Mathjock. A label such as 12.3 means Exercise 3 in Section 12; A.3 means
Exercise 3 in the matrix appendix. Section 0 is included.
These are independently written solutions; brief descriptions
identify the exercises, and the textbook supplies their original statements.\par
\textbf{How to read the solutions.} Identify the assumptions and the conclusion
before beginning a proof. Keep track of the underlying set, its operations,
and the hypotheses of each theorem. For a calculation, follow the intermediate
steps and check the result in the original definition. A counterexample must
satisfy all assumptions and fail the proposed conclusion.\par
}
